\begin{table}[htbp]
\centering
\footnotesize
\caption{Uncertainty Around the Architecture-vs-Information Decomposition}
\label{tab:disentangling_uncertainty}
\begin{tabular}{lccc}
\toprule
Contrast & Estimate & 95\% CI & Excludes 0 \\
\midrule
Architecture average & 477.8 & [319.5, 582.1] & Yes \\
Architecture at Parsimonious & 384.8 & [241.0, 467.1] & Yes \\
Architecture at BoP & 570.9 & [335.5, 729.7] & Yes \\
Information average & 78.1 & [17.9, 111.8] & Yes \\
Information at MS-VAR & -14.9 & [-43.6, 20.5] & No \\
Information at XGBoost & 171.2 & [35.2, 262.2] & Yes \\
Interaction (DiD) & 186.1 & [23.0, 300.8] & Yes \\
\bottomrule
\end{tabular}
\begin{tablenotes}
\footnotesize
\item \textit{Notes:} Positive architecture effects mean XGBoost has higher RMSE than MS-VAR, so larger values favour the regime-switching architecture. Positive information effects mean the BoP specification raises RMSE relative to the parsimonious specification under the paper's sign convention. Confidence intervals use a stationary block bootstrap over rolling-origin dates (2023:01--2025:03; $n=27$, 10 replications, block length 3).
\end{tablenotes}
\end{table}
